\infographicchapter{Cardiac CT}{cardiac_ct.png}
\label{ch:infographic-cardiac-ct}

\section{Overview}

Cardiac computed tomography (cardiac CT) is a quick, noninvasive X-ray scan that takes detailed pictures of your beating heart and its blood vessels. Because the heart is always moving, a regular CT scan would look blurry. To keep the images crisp, a cardiac CT scanner uses small heart-monitor stickers on your chest (an electrocardiogram, or ECG) to match each X-ray picture to your heartbeat. In most scans, the machine takes pictures only during the brief, quiet moment between beats, which keeps the X-ray dose low. In other cases, it scans through the entire heartbeat so doctors can also watch how well the heart pumps.

Cardiac CT is mainly used for two types of exams. The first is a coronary artery calcium (CAC) scan, which uses no contrast dye. It checks for hardened calcium deposits (calcified plaque) inside the heart's arteries and gives you a number called an Agatston score to help estimate your future risk of heart disease. The second is a coronary CT angiogram (CCTA), which uses a special dye injected into a vein in your arm. The dye lights up the blood vessels so doctors can see soft and hardened plaque, narrowed arteries, unusual vessel shapes, bypass grafts, and heart valves (such as when planning a valve replacement). Doctors often order a CCTA for people with chest pain who have a low-to-moderate chance of blocked arteries, because a normal scan is very reliable at ruling out serious blockages.

The biggest benefits of cardiac CT are its clear, detailed views, fast scan time, and ability to rule out blocked arteries without threading a tube into your heart. Its main limits are that a fast or irregular heartbeat can blur the pictures---so you may be given a mild medicine (a beta-blocker) to slow your pulse below 60--65 beats per minute, along with a spray or tablet under your tongue (nitroglycerin) to gently widen the heart's arteries. Heavy calcium buildup can also make narrow spots look larger than they really are. As with any CT scan, your care team takes steps to keep X-ray exposure low and checks for dye allergies and kidney health before using contrast.

\section{Underlying Physical Principles}
\begin{itemize}
    \item \textbf{High-Speed Heart Scanners:} Cardiac CT uses fast-spinning X-ray machines---sometimes with two X-ray tubes rotating at once---that capture each image in a tiny fraction of a second (about 66 to 175 milliseconds). This speed freezes the motion of the beating heart, much like a fast camera shutter freezes a runner.
    \item \textbf{Heartbeat Timing (ECG Gating):} Three small adhesive patches on your chest track your heart's electrical rhythm in real time. With ``prospective'' timing, the scanner turns on the X-ray beam only during the calm resting phase between heartbeats to keep radiation low. With ``retrospective'' timing, it scans through the whole beat so doctors can create a moving movie of the heart.
    \item \textbf{Lighting Up Blood Vessels and Calcium:} During a contrast scan, iodine-based dye flows through the heart's arteries and absorbs X-rays strongly, making the blood vessels stand out brightly. Lowering the X-ray energy level makes the dye shine even more clearly while lowering radiation to the chest. In a calcium scan without dye, hardened calcium spots naturally absorb more X-rays than soft tissue, making them easy to spot and measure.
    \item \textbf{Keeping the Heart Calm and Open:} A gentle heart-slowing medicine (a beta-blocker) gives the heart a longer pause between beats for clearer pictures. A small dose of nitroglycerin relaxes and widens the tiny arteries on the heart's surface so the scanner can capture their fine details.
    \item \textbf{Smart Computer Processing:} Advanced computer programs clean up image noise and reduce bright glare from heavy calcium or metal stents. Using the 3D pictures of your arteries, specialized software can even estimate how well blood flows past a narrowed spot---a test called FFR-CT---without an invasive procedure.
\end{itemize}

\section{Imaging Mechanism}
\begin{enumerate}
    \item \textbf{Setting Up the Heart Monitor and IV:} The technologist places three small ECG stickers on your chest to link your heartbeat to the scanner. For a contrast scan, a nurse or technologist places an IV line in a vein in your arm and checks that your heart rate is calm (usually under 60--65 beats per minute). You may also receive a quick spray or tablet of nitroglycerin under your tongue.
    \item \textbf{Timed Scan During a Short Breath-Hold:} You take a gentle breath and hold it for 10--15 seconds. As the contrast dye reaches your heart, the scanner automatically starts taking pictures timed to your heartbeat---either in short bursts between beats or across the full heart cycle.
    \item \textbf{Capturing Fast X-Ray Snapshots:} As the ring around you spins rapidly, rows of detectors collect X-ray data stamped with the exact moment of your heartbeat, freezing the motion of the tiny coronary arteries.
    \item \textbf{Building Clear 3D Heart Images:} The scanner's computer sorts the snapshots by heartbeat phase, removes blur, and stacks them into razor-thin slices (about half a millimeter thick) or moving clips of your beating heart.
    \item \textbf{Mapping the Arteries and Valves:} On a specialized workstation, the doctor unfolds each twisting coronary branch---the left main, left anterior descending, circumflex, and right coronary arteries---into straightened 2D and 3D views (curved planar reformats), notes which main vessel supplies the back of the heart (coronary dominance), calculates your calcium score if done, and measures blood flow or valve sizes to guide your care.
\end{enumerate}

\section{Key Technical Parameters}
\noindent
\begin{tabularx}{\textwidth}{@{}l X@{}}
\toprule
\textbf{Parameter} & \textbf{Typical Value / Description} \\
\midrule
Physical Principle & Heartbeat-timed (ECG-gated) X-ray CT scanning captured during the heart's quiet resting phase \\
Energy / Signal Source & Fast-rotating CT scanner (64--320 detector rows or dual-source) paired with a 3-lead heart monitor \\
Typical Frequency / Energy & 80--120 kVp X-ray beam adjusted to body size, capturing images in 66--175 milliseconds \\
Spatial Resolution & Very fine detail, with thin image slices about 0.5--0.625 mm thick \\
Typical Scan Duration & 5--15 seconds of actual scanning during a breath-hold (15--30 minutes total in the room) \\
Contrast Agent & None for calcium scoring; iodine-based IV contrast dye (50--90 mL) for coronary CT angiography \\
Ionizing Radiation & Yes (about 0.8--1.5 mSv for calcium scoring; 1--5 mSv for standard CCTA; 5--12 mSv for full-cycle scans) \\
\bottomrule
\end{tabularx}

\section{Common Symptoms That Prompt This Scan}
\begin{itemize}
    \item \textbf{Chest Pain, Tightness, or Shortness of Breath:} New or unexplained chest pressure, squeezing, or breathlessness with exertion when doctors want to check for narrowed heart arteries noninvasively.
    \item \textbf{Unclear Stress Test Results:} Prompted when an exercise stress test or ECG gives borderline results and doctors need a direct 3D look at the coronary arteries.
    \item \textbf{No Symptoms---Heart-Risk Check (Calcium Score):} Frequently ordered in asymptomatic adults with intermediate heart risk (such as family history or high cholesterol) to measure hidden artery calcium.
    \item \textbf{Planning Heart Valve or Rhythm Procedures:} Shortness of breath, fatigue, or irregular heartbeats (atrial fibrillation) when cardiologists need exact 3D heart measurements before valve replacement or ablation.
\end{itemize}

\section{Clinical Usages}

\subsection{Primary Usages}
\begin{itemize}
    \item \textbf{Checking Chest Pain for Blocked Arteries:} Coronary CT angiography (CCTA) is a top noninvasive test for stable or sudden chest pain when the chance of blocked heart arteries is low to moderate. A clear scan strongly rules out dangerous blockages.
    \item \textbf{Estimating Future Heart Disease Risk:} A non-contrast coronary artery calcium (CAC) scan helps measure hidden plaque buildup in people without symptoms, helping you and your doctor decide whether preventive steps like cholesterol-lowering statins are right for you.
    \item \textbf{Planning Heart Valve Replacement:} Detailed heartbeat-timed CT scans measure the exact size of the aortic valve, nearby arteries, and leg blood vessels before minimally invasive valve procedures such as transcatheter aortic valve replacement (TAVR).
    \item \textbf{Checking Unusual Arteries and Bypass Grafts:} Clearly maps heart arteries that formed in unusual places, bulges in vessel walls (aneurysms), and whether past heart bypass grafts remain open and healthy.
    \item \textbf{Guiding Heart Rhythm and Pouch Procedures:} Maps the veins entering the heart before catheter ablation for atrial fibrillation, and checks or measures a small pouch called the left atrial appendage to look for clots or fit a closure device.
\end{itemize}

\subsection{Secondary Usages}
\begin{itemize}
    \item \textbf{Spotting High-Risk Plaque and Measuring Blood Flow:} Helps identify softer, fragile plaques that could trigger a heart attack and supports computer-based blood flow estimates (FFR-CT) to see if a narrowed spot truly limits blood flow.
    \item \textbf{Checking Artificial Heart Valves and Infections:} Detects blood clots, extra tissue growth, leaks, or pockets of infection around artificial heart valves when ultrasound pictures are blocked by metal parts.
    \item \textbf{Evaluating Congenital Heart Conditions and Masses:} Maps complex heart anatomy in adults born with heart defects and checks heart growths or calcium spots when cardiac MRI is not an option.
    \item \textbf{Checking the Heart Sac and Pumping Strength:} Looks for thickening or calcium in the protective sac around the heart (pericardium) and can measure how well the heart's pumping chambers squeeze when other tests are unclear.
\end{itemize}

\section{Advantages}
\begin{itemize}
    \item \textbf{Very Reliable at Ruling Out Blockages:} In people with low-to-moderate risk, a normal coronary CT scan safely rules out serious artery blockages (with a negative predictive value of about 93\%) and helps many patients avoid an invasive heart catheterization. Because an abnormal finding is somewhat less conclusive on its own (roughly 82\% positive predictive value), borderline narrowings are sometimes confirmed with a stress test or catheterization.
    \item \textbf{Shows Both Hard and Soft Plaque:} Unlike tests that only look at blood flow, cardiac CT shows the artery walls directly, revealing early soft plaque, hardened calcium, unusual vessel paths, and bypass grafts.
    \item \textbf{Precise Measurements for Valve Procedures:} Produces crisp 3D measurements of heart valves and blood vessels that are essential for safely sizing and planning modern catheter-based valve treatments.
    \item \textbf{Fast Scan with Long-Term Value:} The actual scan takes only seconds, and calcium scoring or computer flow analysis (FFR-CT) gives your doctor clear information to guide long-term prevention.
    \item \textbf{Checks Nearby Chest Organs Too:} Because the scan covers the chest, doctors can also see the main artery (aorta), lung blood vessels, heart sac, and lungs to spot other possible causes of chest pain.
\end{itemize}

\section{Disadvantages}
\begin{itemize}
    \item \textbf{Sensitive to Fast or Irregular Heartbeats:} A rapid pulse, skipped beats, or atrial fibrillation can blur the pictures or cause steps in the image, so patients often need beta-blockers to slow the heart and nitroglycerin to widen the arteries.
    \item \textbf{Glare from Heavy Calcium or Stents:} Thick calcium deposits and small metal stents can create bright glare (called blooming) on the scan, making a narrowed artery look tighter than it really is.
    \item \textbf{X-Ray Radiation Exposure:} Modern heartbeat-timed scans use much lower doses than older machines, but full-cycle or repeated scans still expose sensitive chest tissues like the breasts and lungs to ionizing radiation.
    \item \textbf{Contrast Dye and Medication Limits:} Contrast scans carry small risks of allergic-like reactions or kidney strain in people with severe kidney disease. Also, beta-blockers or nitroglycerin may not be safe for people with severe asthma, very low blood pressure, or recent use of certain erectile-dysfunction medicines.
    \item \textbf{Diagnosis Only, Not Treatment:} Cardiac CT cannot open a blocked artery during the scan. If a serious or borderline blockage is found, you may still need a stress test, FFR-CT, or an invasive heart catheterization to treat it.
\end{itemize}

\section{Contraindications}

\subsection{Absolute}
\begin{itemize}
    \item \textbf{Recent Use of PDE-5 Inhibitor Medicines (Before Nitroglycerin):} If you have taken sildenafil or vardenafil within the past 24 hours, or tadalafil within the past 48 hours, you must NOT receive nitroglycerin. Combining these medicines can cause a sudden, dangerous, and hard-to-treat drop in blood pressure.
\end{itemize}

\subsection{Relative}
\begin{itemize}
    \item \textbf{Contrast Dye Allergy or Severe Kidney Disease:} If you have had a severe allergic-like reaction to iodine contrast dye or have severely reduced kidney function (eGFR below 30 mL/min/1.73 m\textsuperscript{2}), your doctor may choose a dye-free and radiation-free test such as a stress echocardiogram or stress cardiac MRI instead of an elective CCTA.
    \item \textbf{Conditions That Limit Beta-Blockers:} Uncontrolled heart failure, severe asthma or reactive airway disease, certain serious heart-rhythm blocks, a very slow pulse, or very low blood pressure may prevent the safe use of heart-slowing beta-blockers.
    \item \textbf{Pregnancy or Active Heart Attack:} Routine cardiac CT is usually delayed during pregnancy to avoid radiation. People having an active major heart attack or unstable vital signs need immediate emergency heart catheterization rather than a CT scan.
\end{itemize}

\section{Patient Safety Guidelines}
\begin{itemize}
    \item \textbf{Keeping Radiation as Low as Possible:} Following medical safety guidelines, imaging teams use heartbeat-timed scanning, lower X-ray energy settings matched to your body size, modern noise-reducing software, and focused scan areas to protect your breasts and lungs. Full-cycle scanning is saved only for cases where it is truly needed.
    \item \textbf{Safe Use of Heart-Slowing Medicine:} Before giving an oral or IV beta-blocker, the staff checks for uncontrolled heart failure, severe asthma, heart-rhythm blocks, very slow pulse, or low blood pressure, and watches your heart rhythm and blood pressure closely throughout the visit.
    \item \textbf{Careful Screening Before Nitroglycerin:} Before giving nitroglycerin under your tongue, the team checks that you have not taken sildenafil, vardenafil, or tadalafil in the last 24--48 hours, and screens for severe aortic valve narrowing, certain thickened heart-muscle conditions, or low blood pressure (top number below 90--100 mmHg).
    \item \textbf{Contrast Dye and Kidney Protection:} For contrast scans, staff review any past dye reactions, check kidney blood tests in people at higher risk, make sure the IV line is secure before the fast dye injection, and keep emergency medicines ready at all times. If you have had a past reaction and still need contrast, your doctor may prescribe pre-scan medicines such as a steroid (like prednisone) and an antihistamine (like diphenhydramine). People with kidney disease or diabetes may also receive extra fluids afterward to help flush the iodine dye from the body.
    \item \textbf{Pregnancy and Children's Safety:} Staff check for possible pregnancy before scanning and postpone routine heart CT scans during pregnancy. When children with congenital heart conditions need a scan, teams use ultra-low-dose child-specific settings.
\end{itemize}

\section{Prerequisite Tests}
\begin{itemize}
    \item \textbf{Heart Risk Review and 12-Lead ECG:} Before ordering a cardiac CT, your doctor reviews your symptoms and medical history and usually performs a standard 12-lead ECG. For sudden chest pain in the emergency room, a blood test called high-sensitivity troponin is also checked to rule out an active heart attack.
    \item \textbf{Kidney Blood Test (Only for Contrast Scans):} A calcium score scan without dye needs no lab tests. For a contrast CCTA or valve-planning scan, your doctor may check a recent blood test for kidney function (creatinine and eGFR) if you have kidney disease, diabetes, or other risk factors.
    \item \textbf{Pre-Scan Pulse, Blood Pressure, and Rhythm Check:} Right before a contrast scan, the nurse or technologist measures your resting pulse, blood pressure, and heart rhythm to decide if you need a beta-blocker (aiming for under 60--65 beats per minute) and to confirm your blood pressure is high enough for nitroglycerin.
    \item \textbf{Medication, Allergy, and Pregnancy Check:} The care team asks about any use of erectile-dysfunction or pulmonary-hypertension medicines (like sildenafil, vardenafil, or tadalafil) in the last 24--48 hours, asthma, heart-rhythm issues, past contrast dye reactions, diabetes medicines like metformin, and possible pregnancy.
\end{itemize}

\section{Preparing for the Scan}
\begin{itemize}
    \item \textbf{Skip Caffeine and Eat Lightly Before Contrast Scans:} Avoid all caffeine---including coffee, tea, soda, energy drinks, and chocolate---as well as nicotine and stimulants for 12--24 hours before a contrast CCTA, since they speed up your heart. Avoid solid food for 3--4 hours beforehand, though drinking plain water is encouraged. (A simple calcium scan usually requires no fasting.)
    \item \textbf{Take Regular Medicines and Follow Special Rules:} Keep taking your regular heart medicines and any oral beta-blocker prescribed to calm your pulse before the scan. Strictly avoid sildenafil, vardenafil, or tadalafil for 24--48 hours beforehand so nitroglycerin can be used safely. Ask your doctor if you need to pause metformin or take allergy medicine ahead of time.
    \item \textbf{Dress Simply and Keep Your Chest Clean:} Wear a loose two-piece outfit or plan to change into a hospital gown. Remove necklaces, watches, underwire bras, and metal zippers over the chest, and skip chest lotions or oils that morning so the ECG stickers stay firmly in place.
    \item \textbf{Expect an IV Line and Artery-Widening Spray:} For a contrast scan, a nurse places a sturdy IV line in your arm for the dye (and extra heart-slowing medicine if needed). Just before the pictures are taken, you will usually get a spray or tablet of nitroglycerin under your tongue, which may cause a mild, brief headache or slight lightheadedness.
    \item \textbf{Practice Your Breath-Hold and Plan Ahead:} Practice taking a gentle breath in and holding still for 10--15 seconds without straining. If you weigh more than 300 pounds (136 kilograms), check with the imaging center ahead of time to confirm the scanner table's weight limit. Arrange a ride home if you receive relaxing or allergy medicines, and plan to rest for a few minutes after the scan until your blood pressure feels completely steady.
\end{itemize}

\section{During Scan}
\begin{itemize}
    \item \textbf{Getting Comfortable on the Table:} You will lie flat on your back on a padded CT table with small ECG patches placed on your chest. Both arms rest comfortably above your head on a cushion so they stay out of the picture while the team watches your heart monitor.
    \item \textbf{Following the Breath-Hold Prompt:} A recorded voice from the scanner will ask you to take a gentle breath in and hold completely still for 10--15 seconds. Try not to push or strain, as staying relaxed keeps your heartbeat steady and your pictures sharp.
    \item \textbf{Normal Sensations and Sounds:} If you receive nitroglycerin under your tongue, you may notice a brief tingling taste, mild headache, or facial warmth. As the scanner ring hums and the table glides through the wide opening, the contrast dye injection commonly causes a brief warm flush throughout your body and a mild metallic taste---both fade quickly.
    \item \textbf{Staying in Touch with Your Technologist:} The technologist watches you and your heart rhythm the whole time and can hear you through a two-way intercom. Speak up right away if you feel chest pain, strong dizziness, shortness of breath, itching, hives, or pain or swelling at your IV site.
    \item \textbf{How Long It Takes:} The actual X-ray pictures take only 5--15 seconds (just a few heartbeats). Plan for about 15--30 minutes in the scan room to allow time for the ECG patches, calming your heart rate, and giving nitroglycerin if needed.
\end{itemize}

\section{Normal Appearance}
\begin{itemize}
    \item \textbf{Clean Calcium Score:} A normal coronary artery calcium scan has an Agatston score of 0, meaning no hardened calcium deposits are seen along any of the heart's main arteries.
    \item \textbf{Smooth, Open Heart Arteries:} On a contrast scan, the left and right coronary arteries branch off the aorta in their normal spots and look bright, smooth, and wide open from start to finish, with no soft or hard plaque (reported as CAD-RADS 0).
    \item \textbf{Healthy Heart Chambers and Valves:} The heart's pumping chambers are normal in size and wall thickness, the heart muscle looks even and healthy, the valve flaps are thin and smooth, and the small left atrial pouch fills cleanly with dye without any blood clots.
    \item \textbf{Normal Heart Sac and Major Vessels:} The thin sac around the heart (pericardium) is smooth and under 2 mm thick with no extra fluid or calcium, and the main chest artery (aorta) and lung veins look normal in size and shape.
\end{itemize}

\section{Observing for Abnormal Patterns}
\begin{itemize}
    \item \textbf{Narrowed Arteries and High-Risk Plaque:} Areas where soft, hard, or mixed plaque narrows an artery are graded using a standard scale called CAD-RADS. Doctors also watch for warning signs of fragile plaque---such as outward bulging of the vessel wall, a soft fatty core, tiny flecks of calcium, or a ring-like rim---that carry a higher risk of causing a sudden blockage.
    \item \textbf{Coronary Calcium Buildup:} Bright white spots of calcium along the heart arteries on a non-contrast scan are totaled into the Agatston score. A score of 0 means very low near-term heart risk, while positive scores reflect increasing plaque burden: 1--9 is minimal, 10--99 is mild, 100--299 is moderate, 300--999 is severe, and 1,000 or above indicates extreme plaque buildup.
    \item \textbf{Unusual Vessel Paths and Valve Problems:} The scan can spot a coronary artery taking a risky path squeezed between two large blood vessels, blood clots inside the left atrial pouch (separated from slow-moving blood by a short repeat scan), or abnormal thickening on artificial heart valve leaflets.
    \item \textbf{Heart Muscle and Heart Sac Changes:} Darker patches in the heart muscle wall can show reduced blood flow or an old heart attack (often with thinned muscle). A thickened sac over 4 mm or heavy calcium around the heart can point to a stiff heart sac (constrictive pericarditis).
    \item \textbf{Image Quirks That Can Mimic Disease:} A sudden change in heart rate or an extra beat can cause a ``stair-step'' shift in the image, fast motion can blur the right coronary artery so it looks like soft plaque, and heavy calcium or metal stents can create bright glare that makes a narrowing look worse than it is.
\end{itemize}

\section{Dealing with Uncertainties}
\begin{itemize}
    \item \textbf{Blur from Heartbeat or Breathing:} A fast pulse, irregular rhythm, skipped beats, or moving during the breath-hold can blur parts of the arteries or create step-like offsets, sometimes making a section hard to read clearly.
    \item \textbf{Glare from Heavy Calcium or Stents:} Very dense calcium deposits or small metal stents can shine brightly on the scan and hide the inner blood channel, which can make a mild narrowing look more severe than it truly is.
    \item \textbf{Seeing a Narrowing Versus Knowing Its Effect:} Even on a crystal-clear scan, a moderate narrowing (about 50\%--69\%) does not always tell doctors whether the heart muscle is actually starved of blood during exercise. Also, slow-swirling blood in the left atrial pouch can look like a clot on early pictures.
    \item \textbf{Clearing Up the Picture on the Computer:} Doctors can often solve blurry or unclear spots by rebuilding images from a different moment in the heartbeat, using sharper computer filters for calcium and stents, or taking a brief delayed picture to tell slow blood flow apart from a real clot.
    \item \textbf{Next-Step Tests for Borderline Results:} When a moderate narrowing or heavily calcified area remains uncertain, doctors can use computer-based blood flow analysis (FFR-CT), a stress echocardiogram, a stress cardiac MRI, or a nuclear stress test. Invasive heart catheterization is usually saved for severe or high-risk cases.
\end{itemize}

\section{Augmenting Methods}
\begin{itemize}
    \item \textbf{Computer Blood-Flow Analysis (FFR-CT):} Uses your existing cardiac CT images and advanced computer modeling to estimate whether a moderate narrowing (50\%--69\%) truly reduces blood flow, often preventing unnecessary invasive procedures.
    \item \textbf{Stress Echocardiogram and Nuclear Stress Scans (SPECT/PET):} Check how well blood reaches the heart muscle during exercise or stress medicine when heavy calcium, stents, or borderline narrowings leave questions unanswered.
    \item \textbf{Cardiac MRI:} Uses magnets instead of X-rays to give the clearest measure of heart pumping strength, scar tissue, heart-muscle inflammation (myocarditis), and blood flow under stress.
    \item \textbf{Invasive Heart Catheterization with Tiny Vessel Cameras (IVUS/OCT):} Directly measures tight blockages with a thin pressure wire or tiny camera inside the artery and allows doctors to place a stent right away if needed.
    \item \textbf{Heart Blood Tests and Risk Calculators:} Blood tests like troponin rule out sudden heart injury in the emergency room, while cholesterol panels, lipoprotein(a), and 10-year heart risk calculators work together with your calcium score to guide preventive medications.
\end{itemize}

\section{Comparison with Alternative Modalities}
\begin{itemize}
    \item \textbf{Versus Invasive Heart Catheterization:} Cardiac CT is noninvasive, safer, and ideal for ruling out blockages in low-to-moderate risk chest pain and spotting early plaque. Invasive catheterization is best for active heart attacks or high-risk blockages because doctors can open the artery with a stent during the same procedure.
    \item \textbf{Versus Stress Echocardiogram and Nuclear Stress Testing:} Stress tests measure whether the heart muscle gets enough blood under workload, making them helpful when you already have heavy calcium or stents. Cardiac CT directly shows the arteries, early plaque, unusual vessel paths, and valve sizes for procedures like TAVR.
    \item \textbf{Versus Cardiac MRI:} Cardiac MRI uses no radiation and is the top choice for checking heart muscle scar, inflammation, and pumping volumes. Cardiac CT is much faster and gives far sharper detail of the tiny coronary arteries, calcium deposits, and artificial heart valves.
\end{itemize}

\section{Follow-up Steps}
\begin{itemize}
    \item \textbf{Short Check Before Leaving:} If you received a beta-blocker or nitroglycerin, the team will monitor your pulse and blood pressure until you feel steady and any lightheadedness has passed. They will check your IV site, remind you to drink water to help clear the dye, and review any kidney follow-up if you have severe kidney disease.
    \item \textbf{Clear Results and Prevention Plan:} Your report uses standard CAD-RADS or Agatston calcium scores to help your doctor decide on lifestyle steps, cholesterol or aspirin therapy, or a visit with a cardiologist if significant plaque or narrowing is found.
    \item \textbf{Extra Testing When Needed:} If the scan shows a moderate narrowing (50\%--69\%) or heavy calcium that blocks the view, your doctor may order FFR-CT or a stress test. Severe narrowing in major arteries is usually referred for an invasive heart catheterization.
    \item \textbf{Procedure Planning and Future Scans:} For valve procedures like TAVR or pouch closures, your detailed 3D measurements go straight to the heart team. If your arteries look completely normal and you have no symptoms, you do not need frequent repeat scans; a repeat calcium scan is generally considered no sooner than every 3--5 years if needed.
\end{itemize}

\section{Key References}
\begin{itemize}
    \item Society of Cardiovascular Computed Tomography (SCCT). \textit{SCCT Guidelines for the Performance and Acquisition of Coronary Computed Tomographic Angiography} and \textit{CAD-RADS 2.0 Expert Consensus Document}. J Cardiovasc Comput Tomogr.
    \item Gulati M, Levy PD, Mukherjee D, et al. \textit{AHA/ACC/ASE/CHEST/SAEM/SCCT/SCMR Guideline for the Evaluation and Diagnosis of Chest Pain}. Circulation.
    \item American College of Radiology (ACR). \textit{ACR Appropriateness Criteria: Chest Pain---Possible Acute Coronary Syndrome} and \textit{ACR Manual on Contrast Media}.
    \item Radiological Society of North America (RSNA) and ACR. \textit{RadiologyInfo.org: Coronary CT Angiography (CCTA) and Cardiac CT for Calcium Scoring}.
    \item MedlinePlus Medical Encyclopedia. \textit{Heart CT scan}. U.S. National Library of Medicine, National Institutes of Health (NIH). Available at: \url{https://medlineplus.gov/ency/article/007344.htm}.
    \item Radiopaedia.org. \textit{Cardiac CT}. Radiopaedia. Available at: \url{https://radiopaedia.org/articles/cardiac-ct-1}.
\end{itemize}

\section{Further Reading}
\begin{itemize}
  \item Wikipedia contributors. \textit{Coronary CT angiography}. Wikipedia, The Free Encyclopedia. Available at: \url{https://en.wikipedia.org/wiki/Coronary_CT_angiography}.
\end{itemize}
